\section{Описание реализации}

Программа реализована на языке Python с использованием библиотек
\texttt{NumPy} и \texttt{Matplotlib}. Реализация разделена на две группы:
локальные методы для унимодальных функций и глобальные методы для функций
с несколькими локальными экстремумами.

\subsection{Общая структура программы}

Для унификации результатов работы алгоритмов используется класс
\texttt{OptimizationResult}. Он хранит не только найденную точку минимума,
но и данные, необходимые для последующего анализа:

\begin{itemize}
    \item найденный минимум;
    \item значение функции в найденной точке;
    \item количество вычислений функции;
    \item точки, в которых вычислялась функция;
    \item границы локализующего интервала;
    \item историю сходимости;
    \item историю наилучшего найденного значения для глобальных методов;
    \item текущую оценку константы Липшица.
\end{itemize}

Основные поля результата объединены в структуре данных:

\begin{mintedbox}{python}
@dataclass
class OptimizationResult:
    x_min: float
    f_min: float
    n_evaluations: int
    neighbors: np.ndarray
    evaluation_points: np.ndarray
    n_history: list[int] = field(default_factory=list)
    length_history: list[float] = field(default_factory=list)
    f_best_history: list[float] = field(default_factory=list)
    x_best_history: list[float] = field(default_factory=list)
    lipschitz_M: float | None = None

    @property
    def interval_length(self):
        return self.neighbors[1] - self.neighbors[0]
\end{mintedbox}

Для каждого метода задаются точность по координате минимума
\texttt{eps}, ограничение на число вычислений функции и, при необходимости,
точность по значению функции \texttt{eps\_f}.

\subsection{Локальные методы}

Локальные методы применяются к унимодальным функциям. Их общая идея состоит
в последовательном уменьшении отрезка локализации. В программе реализованы
пассивный поиск, методы дихотомии, Фибоначчи, золотого сечения
и парабол.

\subsubsection{Пассивный поиск}

При пассивном поиске исходный интервал разбивается на равномерную сетку.
Значение функции вычисляется в каждой внутренней точке, после чего выбирается
точка с наименьшим значением функции. Соседние точки сетки образуют итоговый
отрезок локализации.

Число точек выбирается из требования к погрешности координаты:

\begin{equation}
N \geq \frac{b-a}{2\varepsilon}.
\end{equation}

\subsubsection{Метод дихотомии}

На каждой итерации метода дихотомии вычисляются значения функции в двух точках,
расположенных симметрично относительно середины текущего интервала:

\begin{equation}
x_1 = \frac{a+b}{2} - \delta, \qquad
x_2 = \frac{a+b}{2} + \delta.
\end{equation}

Если $f(x_1) < f(x_2)$, правая граница заменяется точкой $x_2$; иначе левая
граница заменяется точкой $x_1$. Процесс продолжается до достижения заданной
длины интервала.

\begin{mintedbox}{python}
    def dichotomy(function):
        #...
        while right - left > eps:
            middle = 0.5 * (left + right)

            x1, x2 = middle - delta, middle + delta
            f1_value, f2_value = function(x1), function(x2)

            if f1_value < f2_value:
                right = x2
            else:
                left = x1
        
        x_min = 0.5 * (left + right)
        #...
\end{mintedbox}

\subsubsection{Методы Фибоначчи и золотого сечения}

Метод Фибоначчи выбирает внутренние точки в соответствии с отношением соседних
чисел Фибоначчи. Перед началом итераций определяется такое число $n$, чтобы
после заданного количества сужений длина интервала стала меньше требуемой
точности. На каждой следующей итерации одна из уже вычисленных точек сохраняется,
поэтому требуется только одно новое вычисление функции.

\begin{mintedbox}{python}
    def fibonacci(function):
        fib_need = (b - a) / (eps - delta)
        n = fibonacci_numbers_until(fib_need)
        # Выбираем в качестве [left, right] нормализованный интервал [0, 1]
        x1 = left + (right - left) * fib(n - 2) / fib(n)
        x2 = left + (right - left) * fib(n - 1) / fib(n)
        #...
        while n > 3:
            if f1_value < f2_value:
                right = x2
                x2, f2_value = x1, f1_value
                n -= 1
                if n == 3: # F(1)/F(3) = F(2)/F(3) = 0.5
                    break
                x1 = left + (right - left) * fib(n - 2) / fib(n)
                f1_value = normalized_function(x1)
                n_evaluations += 1
            else:
                left = x1
                x1, f1_value = x2, f2_value
                n -= 1
                if n == 3:
                    break
                x2 = left + (right - left) * fib(n - 1) / fib(n)
                f2_value = normalized_function(x2)
                n_evaluations += 1
        x2 = x1 + normalized_delta
        f2_value = normalized_function(x2)
        if f1_value < f2_value:
            right = x2
        else:
            left = x1
        # Вычисление в исходных координатах
        x_min = a + (b - a) * (left + right) / 2.0
        #...
\end{mintedbox}

\newpage

Метод золотого сечения использует постоянное отношение

\[
\varphi = \frac{1+\sqrt{5}}{2}.
\]

В отличие от метода Фибоначчи, коэффициент сужения не меняется от итерации к
итерации. Это позволяет применять один и тот же шаг до достижения требуемой
точности. Этот метод можно считать частным случаем метода Фибоначчи при $n \to \infty$.

\begin{mintedbox}{python}
    def golden_section(function):
        #...
        while True:
            if f1_value < f2_value:
                right = x2
                x2, f2_value = x1, f1_value

                if right - left <= normalized_eps:
                    break
                # Критерий остановки по числу вычислений
                elif n_evaluations >= N:
                    break

                x1 = left + (right - left) / phi**2
                f1_value = normalized_function(x1)
                n_evaluations += 1

            else:
                left = x1
                x1, f1_value = x2, f2_value

                if right - left <= normalized_eps:
                    break
                elif n_evaluations >= N:
                    break

                x2 = left + (right - left) / phi
                f2_value = normalized_function(x2)
                n_evaluations += 1

        x_min = a + (b - a) * (left + right) / 2.0
        #...
\end{mintedbox}

\newpage

\subsubsection{Метод парабол}

Метод парабол строит квадратичную аппроксимацию функции по трем точкам
\((x_0, f_0)\), \((x_1, f_1)\) и \((x_2, f_2)\), где
\(f_i = f(x_i)\). Центральная точка должна соответствовать меньшему
значению функции:
\[
x_0 < x_1 < x_2, \qquad f_1 < f_0, \qquad f_1 < f_2.
\]

После чего в качестве новой точки
рассматривает вершину построенной параболы. Для начального выбора тройки точек
используется процедура схожая по принципу с методом дихотомии.

\begin{mintedbox}{python}
    def bracketing(function, interval, N=None):
        #...
        while True:
            #...
            middle, f_middle = (left + right) / 2.0, function(middle)
            n_evaluations += 1

            if f_middle < f_left and f_middle < f_right:
                # Возвращаем искомую тройку точек и значения функции в них,
                # чтобы не пересчитывать повторно.
                return (
                    (left, middle, right),
                    (f_left, f_middle, f_right),
                    n_evaluations,
                    #...
                )
            #...
            if f_left < f_middle:
                right = middle
                f_right = f_middle
            else:
                left = middle
                f_left = f_middle
\end{mintedbox}

Если параболическая аппроксимация вырождается или вершина находится вне текущего
интервала, выполняется запасной шаг деления интервала. Это делает алгоритм
устойчивее для функций, на которых построение параболы временно ненадежно.

\newpage

\begin{mintedbox}{python}
def parabola(function, N=None, show=True):
    #...
    bracket, values, n_evaluations, evaluation_points, boundary_interval = bracketing(function,(a, b), N=N)
    (x0, x1, x2) = bracket
    (y0, y1, y2) = values

    while x2 - x0 > eps:
        extremum = parabola_vertex((x0, x1, x2), (y0, y1, y2))
        use_parabola = extremum is not None

        if use_parabola and (not x0 < extremum < x2 or np.isclose(extremum, x1)):
            use_parabola = False

        if not use_parabola: # Запасной шаг деления интервала
            if x1 - x0 > x2 - x1:
                extremum = 0.5 * (x0 + x1)
            else:
                extremum = 0.5 * (x1 + x2)

        y_extremum = function(extremum)
        n_evaluations += 1

        if extremum < x1:
            if y_extremum < y1:
                x2, y2 = x1, y1
                x1, y1 = extremum, y_extremum
            else:
                x0, y0 = extremum, y_extremum
        else:
            if y_extremum < y1:
                x0, y0 = x1, y1
                x1, y1 = extremum, y_extremum
            else:
                x2, y2 = extremum, y_extremum

        n_history.append(n_evaluations)
        length_history.append(x2 - x0)
    x_min = 0.5 * (x0 + x2)
    #...
\end{mintedbox}

\subsection{Глобальные методы}

Для мультимодальных функций в программе реализованы равномерный перебор
и метод ломаных. В отличие от локальных методов, они не предполагают, что
на всем исходном интервале существует только один минимум.

\subsubsection{Равномерный перебор}

Метод перебора строит равномерную сетку из $N$ точек и вычисляет функцию в каждой
из них. Найденная точка является приближением глобального минимума, а половина
длины шага сетки задает оценку погрешности по координате:

\[
\|x_{\min} - x^*\| \leq \frac{b-a}{2N}.
\]

Для анализа процесса поиска в программе сохраняются последовательности лучших
найденных значений $x_k$ и $f_k$ после каждого вычисления функции.

\subsubsection{Метод ломаных}

Реализация этого метода в работе основана на адаптивном подборе константы Липшица:
на каждой итерации ее значение обновляется по максимальному коэффициенту наклона 
между соседними вычисленными точками.

\begin{equation}
M_{\mathrm{new}} = C_{\mathrm{safety}} \max_i
\left|
\frac{f(x_{i+1})-f(x_i)}{x_{i+1}-x_i}
\right|.
\end{equation}

Множитель $C_{\mathrm{safety}}$ \footnote{Значение множителя 
$C_{\mathrm{safety}} = 1.09$ подобрано
экспериментально.} используется как запас относительно
наблюдаемого максимального наклона.
Без этого запаса нижние ломаные могли бы проходить через уже вычисленные
точки, что снижало бы эффективность локализации минимума.

\begin{mintedbox}{python}
    def update_lipschitz(x_points, y_points, current_M=-np.inf, safety=1.09,):
        order = np.argsort(x_points)
        x_sorted = np.array(x_points)[order]
        y_sorted = np.array(y_points)[order]
        
        slopes = np.abs(np.diff(y_sorted) / np.diff(x_sorted))
        new_estimate = safety * np.max(slopes)

        return max(current_M, new_estimate)
\end{mintedbox}

\newpage

\begin{mintedbox}{python}
def broken_lines(function, N=None, show=True):
    #...
    current_M = max(update_lipschitz(x_points, y_values), 1e-6)
    while True:
        R_values = []
        t_values = []

        for j in range(len(x_points) - 1):
            #... вычисление координат точки пересечения ломаных
            t = 0.5 * (x_left + x_right) - (y_right - y_left) / (2.0 * current_M)
            R = 0.5 * (y_left + y_right) - current_M * (x_right - x_left) / 2.0

            R_values.append(R)
            t_values.append(t)
        # Наименьшее значение R соответствует наиболее перспективной точке
        # для вычисления функции
        best_index = np.argmin(R_values)
        lower_bound = R_values[best_index]

        # Нестрогий критерий остановки
        if np.min(y_values) - lower_bound < eps_f:
            break
        if n_evaluations >= max_evaluations:
            break

        x_new = t_values[best_index]
        y_new = function(x_new)

        if y_new <= min(y_values):
            x_best_history.append(x_new)
        else:
            x_best_history.append(x_best_history[-1])

        insert_index = np.searchsorted(x_points, x_new)
        x_points = np.insert(x_points, insert_index, x_new)
        y_values = np.insert(y_values, insert_index, y_new)
        current_M = update_lipschitz(x_points, y_values, current_M)

    index_min = np.argmin(y_values)
    x_min = x_points[index_min]
\end{mintedbox}

\subsection{Визуализация и сравнение результатов}

Функция \texttt{visualize} строит два графика. На первом показывается функция
на всем исходном интервале, точки вычислений и найденный минимум. На втором
отображается увеличенная окрестность итогового отрезка локализации. Графики
сохраняются в каталоге \\ \texttt{results/<method>/} с именем тестовой функции.

Для исследования сходимости сохраняются число вычислений функции и длина
текущего интервала локализации. Эти данные используются для построения графика
в логарифмическом масштабе. Отдельно реализовано сравнение методов при одинаковом
числе вычислений функции, что позволяет сравнивать не только итоговую точность,
но и вычислительную эффективность алгоритмов.
\newpage
